% TOP_PROBLEM: {"id": 2271, "title": "Erdős Problem #708", "category_id": 1, "category": {"id": 1, "name": "number_theory", "display_name": "Number Theory", "description": "Properties of integers, prime numbers, Diophantine equations.", "slug": "number-theory", "order_index": 1, "created_at": "2026-07-31T15:26:25.670Z"}, "status": "open"}
% FIRST_POSTED: null
% ENTRY_KIND: statement_only
% Imported from: unsolvedmath_additional_500.tex; UnsolvedMath ID 2271
% !TeX program = lualatex
% Compile twice: lualatex 2271.tex
% Self-contained: no external bibliography, images, or \input files.
% Numeric section labels and visible numbers are original UnsolvedMath IDs.
\documentclass[11pt,a4paper]{article}
\usepackage[margin=24mm,headheight=14pt]{geometry}
\usepackage{fontspec}
\setmainfont{Latin Modern Roman}
\setsansfont{Latin Modern Sans}
\setmonofont{Latin Modern Mono}
\usepackage{amsmath,amssymb,mathtools}
\usepackage{unicode-math}
\setmathfont{Latin Modern Math}
\usepackage{microtype}
\usepackage{enumitem,xurl,needspace}
\usepackage[dvipsnames]{xcolor}
\usepackage{hyperref}
\hypersetup{unicode=true,colorlinks=true,linkcolor=MidnightBlue,urlcolor=MidnightBlue,
 pdftitle={UnsolvedMath Problem 2271},pdfauthor={Source-annotated compilation},bookmarksnumbered=true}
\usepackage{bookmark,tocloft,titlesec,fancyhdr}
\setlength{\cftsecnumwidth}{4.2em}
\cftsetpnumwidth{2.5em}
\cftsetrmarg{4em}
\titleformat{\section}{\Large\bfseries\raggedright}{\thesection}{0.7em}{}
\pagestyle{fancy}\fancyhf{}
\fancyhead[L]{\small\sffamily Additional UnsolvedMath statements}
\fancyhead[R]{\small Original catalogue IDs}
\fancyfoot[C]{\thepage}
\setlength{\parindent}{0pt}
\setlength{\parskip}{5pt plus 1pt}
\setlength{\emergencystretch}{3em}
\setcounter{tocdepth}{1}\setcounter{secnumdepth}{1}
\setlist[itemize]{leftmargin=3em,itemsep=4pt,topsep=4pt,parsep=0pt}
\allowdisplaybreaks
\newcommand{\N}{\mathbb{N}}
\newcommand{\Z}{\mathbb{Z}}
\newcommand{\Q}{\mathbb{Q}}
\newcommand{\R}{\mathbb{R}}
\newcommand{\C}{\mathbb{C}}
\newcommand{\F}{\mathbb{F}}
\newcommand{\A}{\mathbb{A}}
\newcommand{\PP}{\mathbb{P}}
\newcommand{\E}{\mathbb{E}}
\newcommand{\eps}{\varepsilon}
\newcommand{\dd}{\,\mathrm d}
\newcommand{\sref}[2]{\hyperlink{source:#1:#2}{[S#2]}}
\newif\ifshowcatalogue
\showcataloguetrue
\newcommand{\cataloguescope}[1]{\ifshowcatalogue
 \paragraph{Statement recorded in the supplied source.}
 #1\fi}
\begin{document}
% UnsolvedMath ID 2271; source catalogue code EP-708

\setcounter{section}{2270}

\section{Covering a Product by Few Factors from Any Long Interval}\label{2271}

{\small\sffamily UnsolvedMath ID 2271 · Number Theory}

\subsection{Definitions and mathematical statement}

\paragraph{Shared notation.}Unless a section says otherwise, $\N=\{1,2,\ldots\}$,
$\Z$ is the ring of integers, $\Q,\R,\C$ are the rational, real, and complex
fields, $[N]=\{1,\ldots,\lfloor N\rfloor\}$, and $\log$ is the natural
logarithm. For real functions of a parameter tending to infinity,
$f\ll g$ and $f=O(g)$ mean $|f|\leq Cg$ eventually for some fixed $C$;
$f\gg g$ reverses this comparison; $f=o(g)$ means $f/g\to0$;
$f\sim g$ means $f/g\to1$; and $f\asymp g$ means both order bounds.
A subscript on $O$ or $\ll$ specifies allowed constant dependence.
The expression $n^{o(1)}$ in an upper bound means growth smaller than
$n^\epsilon$ for each fixed $\epsilon>0$. Local definitions govern any
exception, finite-field convention, or use of the same symbol for another
object. An asymptotic claim is not automatically an effective algorithm.



The intended extremal quantity is an upper bound on the number of selected factors needed. The source writes $|B|=g(n)$, whereas the usual minimization is over $|B|\leq g(n)$; exact-size padding additionally requires a sufficiently long interval. This wording difference is retained in the source statement rather than concealed. \sref{2271}{1} \sref{2271}{2}

\paragraph{Statement recorded in the supplied source.}
Let $g(n)$ be minimal such that for any $A\subseteq [2,\infty)\cap \mathbb{N}$ with $\lvert A\rvert =n$ and any set $I$ of $\max(A)$ consecutive integers there exists some $B\subseteq I$ with $\lvert B\rvert=g(n)$ such that $ \prod_{a\in A} a \mid \prod_{b\in B}b. $ Is it true that $ g(n) \leq (2+o(1))n? $ Or perhaps even $g(n)\leq 2n$? \sref{2271}{1}

\paragraph{Scope and interpretation.}

The record leaves an exact-size versus at-most-size issue. The source statement is preserved and the conventional bound interpretation is identified; this section does not certify the two literal formulations as equivalent in every boundary case. \sref{2271}{1}

\subsection{Short English statement}

Can a product of prescribed integers always divide the product of at most twice as many distinct integers chosen from any sufficiently long interval?

\subsection{Sources}

{\small

\begin{itemize}

\item[\hypertarget{source:2271:1}{S1}] ULAM.AI, \emph{UnsolvedMath}, supplied \texttt{problems.json}, record 2271 (EP-708), statement and background. The embedded literature-review date is 2026-08-17; that is the archive’s review date, not a fresh certification. Dataset landing page: \url{https://huggingface.co/datasets/ulamai/UnsolvedMath}.

\item[\hypertarget{source:2271:2}{S2}] T. F. Bloom, \emph{Erdős Problems}, Problem 708, \url{https://www.erdosproblems.com/708}. Reference imported from the supplied record; the live problem page was not independently checked for this compilation.

\item[\hypertarget{source:2271:3}{S3}] P. Erdős and J. Surányi, Bemerkungen zu einer Aufgabe eines mathematischen Wettbewerbs, Mat. Lapok 10 (1959), 39-48. \url{https://users.renyi.hu/~p_erdos/1959-07.pdf}. Bibliographic information imported from the supplied record.

\item[\hypertarget{source:2271:4}{S4}] Erd\"{o}s, P., Some of my forgotten problems in number theory. Hardy-Ramanujan J. (1992), 34-50. Bibliographic information imported from the supplied record.

\end{itemize}

}

\label{end:2271}
\end{document}
